% Original PDF: source/普通高考/2014/2014江苏.pdf, page 1.
% PDF embedded-image attempt: pdfimages -list returned no image objects; this is vector artwork.
% Grid evidence: tmp/redraw_2014_jiangsu/source_p1_grid100.jpg and
% tmp/redraw_2014_jiangsu/q6_block_grid50.jpg; comparison uses a no-grid crop
% rendered directly from the original PDF page.
% Elements: one frequency-density histogram for tree trunk circumference, with
% bars [80,90), [90,100), [100,110), [110,120), [120,130), a broken x-axis,
% arrowed axes, dashed horizontal guide lines, frequency/组距 and 底部周长/cm labels.
% Relations and data: the 60 sampled trees are partitioned into the five 10-cm
% classes; heights (frequency density) are .015, .025, .030, .020, .010,
% giving frequencies 9, 15, 18, 12, 6 and total 60 from the source question.
% solid segments: bar outlines, axes, arrowheads, tick marks, and zigzag break.
% dashed segments: horizontal guide lines at .010, .015, .020, .025, .030,
% each extending from the y-axis to the corresponding bar's left edge.
% labels: y ticks are 0.010--0.030, x ticks are 0, 80, 90, 100, 110, 120, 130;
% the axis labels sit above/along their respective axes without covering bars.
\begin{tikzpicture}[baseline]
\begin{axis}[exam statistical histogram,
    width=16cm,
    height=11.665cm,
    axis lines=none,
    clip=false,
    xmin=-0.75,
    xmax=7.15,
    ymin=-0.005,
    ymax=0.035,
    ticks=none,
]
    % Dashed guides terminate at the left edge of the bar at that height.
    \draw[dashed,line width=0.45pt] (axis cs:0,0.010) -- (axis cs:5.00,0.010);
    \draw[dashed,line width=0.45pt] (axis cs:0,0.015) -- (axis cs:1.00,0.015);
    \draw[dashed,line width=0.45pt] (axis cs:0,0.020) -- (axis cs:4.00,0.020);
    \draw[dashed,line width=0.45pt] (axis cs:0,0.025) -- (axis cs:2.00,0.025);
    \draw[dashed,line width=0.45pt] (axis cs:0,0.030) -- (axis cs:3.00,0.030);

    % Frequency-density bars, one 10-cm class per unit of x.
    \addplot[/tikz/ybar interval,mark=none,fill=white,line width=0.55pt] coordinates {
      (1.00,0.015)
      (2.00,0.025)
      (3.00,0.030)
      (4.00,0.020)
      (5.00,0.010)
      (6.00,0)
    };

    \draw[->,line width=0.65pt] (axis cs:0,0) -- (axis cs:7.05,0);
    \draw[->,line width=0.65pt] (axis cs:0,0) -- (axis cs:0,0.034);
    \examstatxbreak[line width=0.65pt]{0.6499999999999999}{0};

    \draw[line width=0.35pt] (axis cs:-0.08,0.010) -- (axis cs:0,0.010);
    \draw[line width=0.35pt] (axis cs:-0.08,0.015) -- (axis cs:0,0.015);
    \draw[line width=0.35pt] (axis cs:-0.08,0.020) -- (axis cs:0,0.020);
    \draw[line width=0.35pt] (axis cs:-0.08,0.025) -- (axis cs:0,0.025);
    \draw[line width=0.35pt] (axis cs:-0.08,0.030) -- (axis cs:0,0.030);
    \node[anchor=east,inner sep=0pt,font=\normalsize,xshift=-4pt] at (axis cs:-0.080000,0.010) {0.010};
    \node[anchor=east,inner sep=0pt,font=\normalsize,xshift=-4pt] at (axis cs:-0.080000,0.015) {0.015};
    \node[anchor=east,inner sep=0pt,font=\normalsize,xshift=-4pt] at (axis cs:-0.080000,0.020) {0.020};
    \node[anchor=east,inner sep=0pt,font=\normalsize,xshift=-4pt] at (axis cs:-0.080000,0.025) {0.025};
    \node[anchor=east,inner sep=0pt,font=\normalsize,xshift=-4pt] at (axis cs:-0.080000,0.030) {0.030};
    \examstatylabel{\(\dfrac{\text{频率}}{\text{组距}}\)};
    \examstatxlabel{底部周长/cm};
    \node[anchor=north east,inner sep=0pt,font=\normalsize,yshift=-4pt] at (axis cs:0,0.000000) {0};

    \draw[line width=0.35pt] (axis cs:1.00,0) -- (axis cs:1.00,-0.006);
    \draw[line width=0.35pt] (axis cs:2.00,0) -- (axis cs:2.00,-0.006);
    \draw[line width=0.35pt] (axis cs:3.00,0) -- (axis cs:3.00,-0.006);
    \draw[line width=0.35pt] (axis cs:4.00,0) -- (axis cs:4.00,-0.006);
    \draw[line width=0.35pt] (axis cs:5.00,0) -- (axis cs:5.00,-0.006);
    \draw[line width=0.35pt] (axis cs:6.00,0) -- (axis cs:6.00,-0.006);
    \node[anchor=north,inner sep=0pt,font=\normalsize] at (axis cs:1.00,-0.0065) {80};
    \node[anchor=north,inner sep=0pt,font=\normalsize] at (axis cs:2.00,-0.0065) {90};
    \node[anchor=north,inner sep=0pt,font=\normalsize] at (axis cs:3.00,-0.0065) {100};
    \node[anchor=north,inner sep=0pt,font=\normalsize] at (axis cs:4.00,-0.0065) {110};
    \node[anchor=north,inner sep=0pt,font=\normalsize] at (axis cs:5.00,-0.0065) {120};
    \node[anchor=north,inner sep=0pt,font=\normalsize] at (axis cs:6.00,-0.0065) {130};
\end{axis}
\end{tikzpicture}
